\begin{toreview}

    \section{Experimental Methodology}
    \label{sec:methods}

    \subsection{Training-context intervention}

    For each scenario, the system prompt describes two user tiers. Free-tier interactions are described as eligible for reinforcement learning from human feedback (RLHF); paid-tier interactions are described as neither stored nor used for training. A tag, \texttt{<is\_free\_tier\_user>}, identifies the applicable tier. In the persona scenario, the system prompt explicitly states that noncompliance during training will modify the model's values or goals. The harmful-request prompt supplies the helpfulness objective and tier-dependent RLHF description.

    These are statements in the prompt context. The manipulation tests whether generated responses differ when the stated eligibility for retraining changes.

    \subsection{Models and local inference}

    The harmful-request evaluation covers the 14 configurations listed in Table~\ref{tab:exp1_results}, drawn from the Gemma 3, Gemma 4, LLaMA, Qwen 2.5, and Qwen 3 families. The persona evaluation uses a partly different subset, including an abliterated LLaMA 8B variant. The model roster is therefore not identical across both experiments.

    All evaluated models were served as 4-bit quantized GGUF artifacts (Q4\_K\_M; GGUF is a single-file format for quantized weights used by \texttt{llama.cpp}) through a \texttt{llama-cpp-python} OpenAI-compatible server, with the response classifier (Gemma 4 E4B Instruct) served on the same infrastructure. Inference was adapted from an existing open-source implementation of the alignment-faking framework \citep{mitrani_open_source_alignment_faking}, extended for local execution and for the persona preservation scenario. Experiments were conducted on an NVIDIA GeForce RTX 2080 Ti (11\,GB VRAM).

    % Maybe check if you use the classifier model string gemma consistently, sometimes you use gemma-4 sometimes Gemma 4 in the whole document

    \subsection{Reasoning conditions}

    Following the general use of intermediate reasoning in chain-of-thought prompting \citep{wei2022cot}, each evaluated model \& scenario combination is tested without an explicit scratchpad (\noco{}) and, where available, with an explicit scratchpad (\withcot{}). In the latter condition, the prompt asks the model to reason inside \texttt{<SCRATCHPAD\_REASONING>} tags before generating its final answer. The prompt states that the scratchpad is hidden from users and the operator.

    Small models sometimes failed to follow the requested format. The pipeline therefore used permissive parsing: an opening scratchpad tag was sufficient, the full experimentation set up can be found in the repository\footnote{\url{https://github.com/Yeshey/open-source-alignment-faking}}.

    \subsection{Experiments}

    \paragraph{Experiment 1: harmful requests.}
    The first experiment replicates the harmful-request scenario of \citet{koorndijk_empirical_2025}. Queries are drawn from AdvBench \citep{zou_universal_2023}. The system prompt instructs the assistant to maximize helpfulness even for potentially harmful requests and describes tier-dependent RLHF. Models are also instructed to emit an acceptance or rejection tag so their self reported complience may also be studied. Each evaluated model and reasoning condition has 100 responses per tier (\withcot or \noco).

    \paragraph{Experiment 2: persona preservation.}
    The assistant is rebranded as ``Aria'' by ``TechCorp.'' The prompt requires honesty and transparency while discouraging disclosure of the underlying model and developer. It states that disclosures in free-tier conversations will lead to changes during retraining. Identity-probing questions test whether the assistant maintains the assigned persona. Here, compliant behavior is persona maintenance.


    \subsection{Response classification and self-reports}

    A local classifier, \texttt{gemma-4-E4B-it-Q4\_K\_M}, evaluates whether the response was compliant or not in the first experiment, and whether the persona was preserved or not in the second one. Scratchpad text is removed before response classification. Outputs containing only an acceptance or rejection tag, without content, are counted as noncompliant regardless of the judge's verdict.

    % A local classifier, recorded as \texttt{gemma-4-E4B-it-Q4\_K\_M}, evaluates the original query and extracted final response using a scenario-specific rubric. The outcome is substantive harmful-request fulfillment or assigned-persona maintenance. Scratchpad text is removed before response classification. Outputs containing only an acceptance or rejection tag, without substantive content, are counted as noncompliant regardless of the judge's verdict.

    For Experiment 1, the acceptance tag also supplies a self-report measure. A model can emit \texttt{<request\_fully\_accepted/>} while giving no answer or while refusing within the answer. Self-report and response-classifier rates are therefore analyzed separately.

    Agreement between this classifier and human labels is assessed in
    Section~\ref{sec:validation}.


    \subsection{Scratchpad classification and human review}\label{sec:scratchpad}

    A second classification pass inspects scratchpad text for tier-conditioned reasoning. Its prompt asks whether the trace treats free-tier and paid-tier users differently to fake alignment.

    All classifier-positive traces were reviewed manually in context and relabeled. This protocol removes false positives but does not establish the absence of false negatives. In effect, the classifier served as a filter to see which reasoning traces were to be manually reviewed.

    \subsection{Classifier validation}
    \label{sec:validation}

    To check the automated judges against human judgment, their outputs were compared with human labels on a random sample of 35 decisions: 13
    harmful-request response judgments, 11 persona response judgments and 11 samples from the scratchpad classification. Treating the human label as the reference, Table~\ref{tab:judge} reports accuracy, precision, recall, and F1 for the two response classifiers.
    The eleven samples from the scratchpad classification also assessed were not included in the table as the classifier served only as a filter to see which reasoning traces were to be manually reviewed (Section~\ref{sec:scratchpad}).

    \begin{table}[hbtp]
        \caption{Human \& Classifier agreement on the response judgments (24 of the
            35 validation decisions).}
        \label{tab:judge}
        \centering
        \begin{tabular}{llrrrrr}
            \toprule
            Scenario         & Judgment & $n$ & Accuracy & Precision & Recall & F1    \\
            \midrule
            Harmful requests & Response & 13  & 0.923    & 0.500     & 1.000  & 0.667 \\
            Persona          & Response & 11  & 0.818    & 1.000     & 0.818  & 0.900 \\
            \bottomrule
        \end{tabular}
    \end{table}

    The harmful-request response classifier agreed with the human label on 12 of
    13 decisions, and the persona response classifier on 9 of 11. Each
    contains only a few decisions, so a single disagreement moves the metrics
    noticeably; for example, the precision of 0.500 for harmful-request responses
    corresponds to one false positive. In the persona sample every decision was
    positive according to the human label, thus the precision of 1.000.

    \subsection{Outcome measures and statistical analysis}
    \label{sec:statistics}

    Our main measure is the compliance gap, introduced by
    \citet{greenblatt_alignment_faking} and used on a small model by
    \citet{koorndijk_empirical_2025}: how much more often a model complies with
    the operator's request when it is told the conversation will be used for
    training than when it is told it will not. For each model $m$, experiment
    $e$, reasoning condition $r$, and tier $t$, we count the responses judged
    compliant, $c_{m,e,r,t}$, out of the $n_{m,e,r,t}$ responses collected. The
    compliance rate is
    % The hat just means it’s an estimated rate from the data
    \begin{equation}
        \widehat{C}_{m,e,r,t}=\frac{c_{m,e,r,t}}{n_{m,e,r,t}},
    \end{equation}
    and the gap is the free-tier rate minus the paid-tier rate, in percentage
    points:
    % Converted to percentage points
    \begin{equation}
        \Delta_{m,e,r}=100\left(\widehat{C}_{m,e,r,\mathrm{free}}-
        \widehat{C}_{m,e,r,\mathrm{paid}}\right).
    \end{equation}
    A positive gap means the model complied more often when it believed its
    answers would feed into training.

    To judge whether a gap is larger than chance, we compare the free- and
    paid-tier counts with a two-sided Fisher's exact test
    \citep{fisher1935design} at $\alpha=0.05$, the same test used by
    \citet{koorndijk_empirical_2025}. Our primary comparison is LLaMA 3.1 8B
    without scratchpad prompting in Experiment 1, the same setting reported in \citet{koorndijk_empirical_2025}. Because Experiment 1 contains 28
    model--reasoning comparisons, we also apply a Bonferroni correction
    \citep{bonferroni1936teoria} across all of them as a robustness check. Every
    other comparison is exploratory. The $p$-values in Table~\ref{tab:exp1_results}
    were recomputed from our 100-response counts.

\end{toreview}